\documentclass[12pt]{scrartcl}

\input{../styles/Packages}
\input{../styles/FormatAndHeader}

\usepackage{xcolor} % Für die Farben im Code

% Definiere das Aussehen des Codes
\lstset{
    language=XML,
    basicstyle=\ttfamily\small,
    numbers=left,           % Zeilennummern links
    stepnumber=1,
    showstringspaces=false, % Keine komischen Unterstriche in Strings
    tabsize=4,
    breaklines=true,        % Automatischer Zeilenumbruch
    breakatwhitespace=false,
    frame=single,           % Ein Rahmen um den Code
    extendedchars=true,
    literate={ö}{{\"o}}1 {ä}{{\"a}}1 {ü}{{\"u}}1 {ß}{{\ss}}1 {Ö}{{\"O}}1 {Ä}{{\"A}}1 {Ü}{{\"U}}1
}

% \stepcounter{countername}: Increases counter
\setcounter{sheetnr}{7} % Number of sheet


\setcounter{exnum}{1} % Exercise number

\begin{document}

    %TODO: Aufgabe 7.1


    %TODO: Aufgabe 7.2


    \exercise{7.3: XML-Instanz}

% Hinweise zur Lösung:
% 1. xs:date erfordert das Format YYYY-MM-DD.
% 2. xs:sequence erzwingt die exakte Reihenfolge der Elemente (Eintrittsdatum MUSS vor Position stehen).
% 3. Das "Gehalt" ist im XSD-Schema nicht definiert und muss weggelassen werden,
%    da das Dokument sonst nicht gültig (valid) gegenüber dem Schema wäre.

\begin{lstlisting}[language=XML]
<?xml version="1.0" encoding="UTF-8"?>
<Personalverwaltung>
  <MitarbeiterInfo>
    <PersonalID>MA-8821</PersonalID>
    <Nachname>Mustermann</Nachname>
    <Eintrittsdatum>2026-06-01</Eintrittsdatum>
    <Position>Senior</Position>
  </MitarbeiterInfo>

  <MitarbeiterInfo>
    <PersonalID>MA-1024</PersonalID>
    <Nachname>Schmidt</Nachname>
    <Eintrittsdatum>2026-01-15</Eintrittsdatum>
    <Position>Manager</Position>
  </MitarbeiterInfo>

  <MitarbeiterInfo>
    <PersonalID>MA-5566</PersonalID>
    <Nachname>Weber</Nachname>
    <Eintrittsdatum>2026-03-20</Eintrittsdatum>
    <Position>Junior</Position>
  </MitarbeiterInfo>
</Personalverwaltung>
\end{lstlisting}


\end{document}
